The GZKST abstraction isolates the seed-hiding invariant common to
seed-tree-based post-quantum signatures, exhibits the published fault attacks
as violations of it, and turns the fault surface into an enumerable grid from
which we derived generic key-recovery procedures with bounded query
complexity. \emph{We demonstrated three clock-glitch attacks against the MEDS}
reference implementation on an ARM Cortex-M4 and evaluated countermeasures
whose cost is small relative to the matrix arithmetic dominating signing.

\chg{Three limitations remain. The glitches targeted the path-construction
function in isolation, not inside a complete signature, so locating the target
in a full signing trace is not yet demonstrated. Model~3 required building the
target function at \texttt{-O0}, since optimised builds remove the targeted
assignment. And the double glitch that defeats verify-before-release has not
been measured.}

%Two directions follow. The first is to quantify the double-glitch attack of
%Section~\ref{sec:counter}, turning a structural observation about shared code
%into a concrete attacker cost. The second is to instantiate the abstraction
%against a scheme still under consideration --- FAEST, SDitH and MQOM all build
%an all-but-one vector commitment from the same tree --- since the empty cells
%of Table~\ref{tab:surfaces} are as much a programme as a classification.